\begin{table}[H]
%\begin{table}[ht]
\caption{History}
\label{tab:arch_history03}
\centering
\begin{tabularx}{\textwidth}{|X |X |X |}
  \hline
  \multicolumn{3}{|l|}{History} \\
  \hline
  \multicolumn{1}{|l}{Target Spec.} & \multicolumn{1}{l}{Current version : 1.1, 2026-11-18} & \multicolumn{1}{l|}{Author: A. Siggelkow}\\
  \multicolumn{1}{|l}{ } & \multicolumn{1}{l}{Previous version : 1.0, 2025-02-22} & \multicolumn{1}{l|}{ }\\
  \hline
  \hline
  Paragraph (in previous version) & Paragraph (in current version) / date & Subjects (major changes since last revision) \\
  \hline
  \hline
  \multicolumn{3}{|l|}{Current Chapter Version: C0.1; Date: 2025-02-22} \\
  \hline
  & 2025-02-22 & First draft (AS) \\
  \hline
  \multicolumn{3}{|l|}{Current Chapter Version: C0.2; Date: 2026-11-18} \\
  \hline
  & \ref{sec:archirq} / 2026-11-18 & Interrupt concept of the peripherals: register structure (different IRQs, IRQs of the same type), interrupt flow, generic register description; from HowTo PerIRQ (AS, Claude) \\
  \hline
\end{tabularx}
\end{table}
